\section{Interaction Classification}
\label{sec:interaction-classification}

The geometric structures developed in previous sections—semantic excitation,
geodesic trajectories, spectral bundles, and holonomy return operators—enable a
formal classification of NDH interaction behavior. This section defines the
interaction classes arising from modal alignment, spectral deformation,
recursion transitions, and holonomy accumulation.

\subsection{Classification Framework}

An interaction is defined by the tuple
\[
\mathcal{I} = (x, \mathbf{s}, \gamma, \mathcal{H}, \mathcal{H}_S, \mathcal{H}_R),
\]
where $x$ is the initial interaction state, $\mathbf{s}$ is the excitation
vector, $\gamma$ is the induced geodesic trajectory, and
$(\mathcal{H}, \mathcal{H}_S, \mathcal{H}_R)$ are the holonomy operators
associated with tangent, spectral, and recursion transport along $\gamma$.
Classification is determined by the structural properties of this tuple.

\subsection{Modal-Aligned Interactions}

An interaction is \emph{modal-aligned} if the geodesic $\gamma$ remains within a
single modal region $\mathcal{M}_i$ and if the holonomy return operator is
trivial:
\[
\mathcal{H}(\gamma) = \mathrm{id}.
\]
Modal-aligned interactions exhibit stable recursion behavior governed by the
operator $R_i$ and minimal spectral deformation. These interactions correspond
to structurally coherent interrogation within a single modal domain.

\subsection{Spectral-Aligned Interactions}

An interaction is \emph{spectral-aligned} if modal-boundary crossings induce
nontrivial spectral deformation while tangent holonomy remains trivial:
\[
\mathcal{H}(\gamma) = \mathrm{id}, \qquad
\mathcal{H}_S(\gamma) \neq \mathrm{id}.
\]
Spectral-aligned interactions arise when semantic excitation induces twisting of
the spectral bundle without altering the underlying modal or recursion
structure. These interactions correspond to spectral drift or rotation driven by
boundary-induced fiber morphisms.

\subsection{Recursion-Aligned Interactions}

An interaction is \emph{recursion-aligned} if the recursion holonomy operator is
nontrivial:
\[
\mathcal{H}_R(\gamma) \neq \mathrm{id}.
\]
These interactions involve cyclic transitions between recursion operators
$(R_A, R_B, R_C)$ along the trajectory. Recursion-aligned interactions correspond
to interrogation paths that traverse modal boundaries in a manner that induces
non-dual recursion effects.

\subsection{Holonomy-Deviant Interactions}

An interaction is \emph{holonomy-deviant} if the tangent holonomy operator is
nontrivial:
\[
\mathcal{H}(\gamma) \neq \mathrm{id}.
\]
Holonomy-deviant interactions exhibit geometric deformation of the tangent
structure, indicating curvature-dependent effects in the cognitive manifold.
These interactions correspond to interrogation loops that accumulate geometric
distortion independent of spectral or recursion behavior.

\subsection{Mixed Interactions}

An interaction is \emph{mixed} if multiple holonomy operators are nontrivial:
\[
\mathcal{H}(\gamma) \neq \mathrm{id}, \qquad
\mathcal{H}_S(\gamma) \neq \mathrm{id}, \qquad
\mathcal{H}_R(\gamma) \neq \mathrm{id}.
\]
Mixed interactions arise from complex interrogation patterns involving modal
transitions, spectral twisting, and recursion cycling. These interactions
represent the full geometric richness of NDH's interaction behavior.

\subsection{Transformer-Flat Interactions}

For comparison, transformer-based systems exhibit \emph{flat} interaction
behavior. In the absence of curvature, modal boundaries, spectral bundles, and
recursion kernels, transformer interactions satisfy:
\[
\mathcal{H}(\gamma) = \mathrm{id}, \qquad
\mathcal{H}_S(\gamma) = \mathrm{id}, \qquad
\mathcal{H}_R(\gamma) = \mathrm{id}.
\]
Transformer-flat interactions lack geometric structure and therefore cannot be
classified using the framework developed here.

\subsection{Summary}

NDH interaction behavior is classified according to modal alignment, spectral
deformation, recursion transitions, and holonomy accumulation. These classes
provide a geometric taxonomy of interrogation, distinguishing NDH from
transformer-based systems and establishing interaction geometry as a coherent
mathematical framework.
